\begin{lemma}{Gordon'sche Paare auf dichten Teilmengen}
             {GordonschePaareAufDichtenTeilmengen}
Seien $c\in\R$ und $(X,Y)$ das Gordon'sche Paar zum Tupel
$(\gamma,\delta,g,c)$.\\
Seien $\emptyset \neq T \subseteq \R^m$ beschr�nkt und
$\emptyset \neq V \subseteq \R^N$ beschr�nkt, sowie $T' \subseteq T$
dicht in $T$ und $V' \subseteq V$ dicht in V.\\
Dann sind $T'\neq \emptyset$, $V' \neq \emptyset$ und es gilt:
\begin{enumerate}[(1)]
	\item \label{GordonschePaareAufDichtenTeilmengen1}
	$ \underset{v \in V}\sup\, \bprod v \gamma _N 
	 =\underset{v \in V'}\sup\,\bprod v \gamma_N$.
	\item \label{GordonschePaareAufDichtenTeilmengen2}
	 $\underset{t \in T}\inf\, c\bprod t \delta_m 
	 =\underset{t \in T'}\inf\, c\bprod t \delta_m\,$ und $\, 
	 \underset{t \in T} \sup\, c\bprod t \delta_m 
	  =\underset{t \in T'}\sup\, c\bprod t \delta_m$.\\
	  Insbesondere gilt
 \[\inf_{t \in T}\sup_{v \in V}X(t, v) = \inf_{t\in T'}\sup_{v \in V'}X(t,v)\, \text{ und }
 \,\sup_{t\in T}\sup_{v\in V}X(t,v)=\sup_{t\in T'}\sup_{v\in V'}X(t,v)\text.\]
	\item \label{GordonschePaareAufDichtenTeilmengen3}
	$\Inf{t \in T}\Sup{v \in V}Y(t, v) = \Inf{t\in T'}\Sup{v \in V}Y(t, v)\,$ und
	$\,\Sup{t \in T}\Sup{v \in V}Y(t, v) = \Sup{t\in T'}\Sup{v \in V}Y(t, v)$.
	\item \label{GordonschePaareAufDichtenTeilmengen4}
	$\Inf{t \in T}\Sup{v \in V}Y(t, v) = \Inf{t\in T'}\Sup{v \in V'}Y(t,v)\,$ und
	$\,\Sup{t \in T}\Sup{v \in V}Y(t, v) = \Sup{t\in T'}\Sup{v \in V'}Y(t, v)$.
\end{enumerate}
\end{lemma}
\begin{proof}
Da die leere Menge stets abgeschlossen ist und 
$\overline {T'}^T = T \neq \emptyset$, sowie 
$\overline{V'}^V= V \neq \emptyset$ gilt, sind $T'$ und $V'$ nichtleer.\\
Da $V$ beschr�nkt ist, gibt es $d \in \R_{>0}$ mit 
$\forall\,v \in V: \norm v_{2, N} \le d$.\\
Sei $\omega \in \Omega$.\\
(1) Die Abbildung $\sprod \cdot {\gamma(\omega)}_N$ ist lipschitzstetig und da
lipschitzstetige Bilder beschr�nkter Mengen beschr�nkt sind, ist die Menge
$\sprod \cdot {\gamma(\omega)}_N(V)$ beschr�nkt.\\
Wegen $V \neq \emptyset$ und $V'$ dicht in $V$, gilt:
\begin{align*}
\left(\sup_{v \in V}\bprod v \gamma_N\right)(\omega)
&= \sup_{v \in V}\sprod v {\gamma(\omega)}_N 
= \sup \left(\sprod \cdot {\gamma(\omega)}_N \right)(V) \\
&\overset{\refclaim{StetigeFunktionenUndDichteTeilmengen}
{StetigeFunktionenUndDichteTeilmengen1}}=
 \sup \left(\sprod \cdot {\gamma(\omega)}_N \right)(V')
 = \left(\sup_{v \in V'}\bprod v \gamma_N\right)(\omega)
\end{align*}
(2) Die Abbildung $c \sprod \cdot {\delta(\omega)}_m$ ist ebenfalls
lipschitzstetig und daher auch die Menge 
$(c \sprod \cdot {\delta(\omega)}_m)(T)$ nach oben und unten beschr�nkt.\\
Da $T \neq \emptyset$ und $T'$ dicht in $T$, folgt wie in (1):
\begin{align*}
\left(\inf_{t \in T}\bprod t \delta_m\right) (\omega)
&\overset{\refclaim{StetigeFunktionenUndDichteTeilmengen}
{StetigeFunktionenUndDichteTeilmengen2}}=
\left(\inf_{t \in T'}\bprod t \delta_m\right) (\omega)\text,\\
\left(\sup_{t \in T}\bprod t \delta_m\right) (\omega)
&\overset{\refclaim{StetigeFunktionenUndDichteTeilmengen}
{StetigeFunktionenUndDichteTeilmengen1}}=
= \left(\sup_{t \in T'}\bprod t \delta_m\right) (\omega)\text.
\end{align*}
Weiter gilt wegen der Beschr�nktheit von $T, T', V, V'$:
\begin{align*}
\inf_{t \in T'}\sup_{v \in V'}X(t, v) &
\overset{\refclaim{BeschraenktheitGordonscherPaare}
{BeschraenktheitGordonscherPaare3}}=
\inf_{t \in T'}c\bprod t \delta_m + \sup_{v \in V'}\bprod v \gamma_N\\
&\overset{(1),(2)}= 
\inf_{t \in T}c\bprod t \delta_m + \sup_{v \in V}\bprod v \gamma_N
\overset{\refclaim{BeschraenktheitGordonscherPaare}
{BeschraenktheitGordonscherPaare3}}=\inf_{t \in T}\sup_{v \in V}X(t, v)\text,\\
\sup_{t \in T'}\sup_{v \in V'}X(t, v) &
\overset{\refclaim{BeschraenktheitGordonscherPaare}
{BeschraenktheitGordonscherPaare4}}=
\sup_{t \in T'}c\bprod t \delta_m + \sup_{v \in V'}\bprod v \gamma_N\\
&\overset{(1),(2)}= 
\sup_{t \in T}c\bprod t \delta_m + \sup_{v \in V}\bprod v \gamma_N
\overset{\refclaim{BeschraenktheitGordonscherPaare}
{BeschraenktheitGordonscherPaare4}}=\sup_{t \in T}\sup_{v \in V}X(t, v)\text.
\end{align*}
(3)/(4) F�r jedes $w\in V$ ist die Abbildung $\sprod \cdot{g(\omega)^\top w}_m$
lipschitzstetig und linear und es gilt f�r alle $x \in \R^m$:
\begin{align}
\abs{\sprod x {g(\omega)^\top w}_m} &\le
\norm x_{2, m}\norm{g(\omega)^\top w}_{2, m} 
\le \norm x_{2, m} \opnorm{g(\omega)^\top} {\R^N}{\R^m} \norm w _{2, N}\notag\\
&\le \norm x_{2, m} \opnorm{g(\omega)^\top} {\R^N}{\R^m} d\text.
\label{gpadteq:1}
\end{align}
Somit ist $\menge{\sprod x {g(\omega)^\top v}}{v \in V}$ beschr�nkt
und $\left(\sprod \cdot {g(\omega)^\top v}_m \right)_{v \in V}$ eine
Familie gleichm��ig lipschitzstetiger Abbildungen.\\
Somit ist nach \ref{GleichmaessigeLipschitzstetigkeit} die Abbildung
\[f: \R^m \rightarrow \R, t \mapsto 
 \sup\menge{\sprod t {g(\omega)^\top v}_m}{v \in V}\]
wohldefiniert und lipschitzstetig.\\
Da $T$ beschr�nkt ist, ist also auch $f(T)$ beschr�nkt, sodass gilt
\begin{align*}
\left(\inf_{t \in T}\sup_{v \in V}Y(t, v) \right)(\omega)&=
\inf_{t \in T} \sup_{v \in V}\sprod {g(\omega)t} v_N =
\inf_{t \in T} \sup_{v \in V}\sprod t {g(\omega)^\top v}_N\\
&= \inf_{t \in T} f(t)
\overset{\refclaim{StetigeFunktionenUndDichteTeilmengen}
{StetigeFunktionenUndDichteTeilmengen2}}=
\inf_{t \in T'} f(t)
 = \left(\inf_{t \in T'}\sup_{v \in V}Y(t, v) \right)(\omega)\text.
\end{align*}
Es ist $T \times V'$ dicht in $T \times V$, wenn man $\R^m \times \R^N$ mit der
Produktnorm ausstattet.\\
Die Abbildung $\sprod{g(\omega)\cdot}{\cdot\cdot}_N$ ist
als Komposition stetiger Funktionen stetig.\\
Daher gilt
\begin{align*}
\left(\sup_{t \in T}\sup_{v \in V}Y(t, v) \right)(\omega)&= 
\sup_{(t,v)\in T\times V} Y(t, v)(\omega) 
= \sup_{(t,v)\in T\times V} \sprod{g(\omega)t}v_N\\
&\overset{\refclaim{StetigeFunktionenUndDichteTeilmengen}
{StetigeFunktionenUndDichteTeilmengen1}}=
\sup_{(t, v) \in T \times V'}\sprod{g(\omega)t}v_N
= \left(\sup_{t \in T}\sup_{v \in V'}Y(t, v) \right)(\omega)
\end{align*}
Damit ist zun�chst (3) gezeigt.\\
F�r jedes $s \in T$ ist die Abbildung $\sprod {g(\omega)s}\cdot_N$
lipschitzstetig, sodass gilt:
\begin{align}
\sup_{v \in V}\sprod{g(\omega)s}v_N
\overset{\refclaim{StetigeFunktionenUndDichteTeilmengen}
{StetigeFunktionenUndDichteTeilmengen2}}= 
\sup_{v \in V'}\sprod{g(\omega)s}v_N\text. \label{gpadteq:2}
\end{align}
Wegen $\eqref{gpadteq:1}$ ist (analog zu $f$) auch 
\[h : \R^m \rightarrow \R, 
t \mapsto \sup\menge{\sprod t {g(\omega)^\top v}_m}{v \in V'}\]
nach \ref{GleichmaessigeLipschitzstetigkeit} wohldefiniert und
lipschitzstetig.\\
Erhalte daher
\begin{align*}
\left(\inf_{t \in T}\sup_{v \in V}Y(t, v) \right)(\omega)&=
\inf_{t \in T} \sup_{v \in V}\sprod {g(\omega)t} v_N
\overset{\eqref{gpadteq:2}}=
\inf_{t \in T} \sup_{v \in V'}\sprod {g(\omega)t} v_N\\
&= \inf_{t \in T} \sup_{v \in V'}\sprod t {g(\omega)^\top v}_N
= \inf_{t \in T} h(t)\overset{\refclaim{StetigeFunktionenUndDichteTeilmengen}
{StetigeFunktionenUndDichteTeilmengen2} }= \inf_{t \in T'} h(t)\\
&= \inf_{t \in T'} \sup_{v \in V'}\sprod t {g(\omega)^\top v}_N
= \left(\inf_{t \in T'}\sup_{v \in V'}Y(t, v) \right)(\omega)\text.
\end{align*}
Es ist auch $T' \times V'$ dicht in $T \times V$, wenn man $\R^m \times \R^N$
mit der Produktnorm ausstattet. Daher gilt auch:
\begin{align*}
\left(\sup_{t \in T}\sup_{v \in V}Y(t, v) \right)(\omega) &=
\sup_{(t, v) \in T\times V}Y(t, v)(\omega) = 
\sup_{(t, v) \in T\times V}\sprod {g(\omega)t} v_N\\
&\overset{\refclaim{StetigeFunktionenUndDichteTeilmengen}
{StetigeFunktionenUndDichteTeilmengen1}}=
\sup_{(t, v) \in T'\times V'}\sprod {g(\omega)t} v_N
= \left(\sup_{t \in T'}\sup_{v \in V'}Y(t, v) \right)(\omega)\text.
\end{align*}
\end{proof}